\input{../common/preamble.tex}

\title{Aarhus Bryghus - Bestilling af fadøl}
\date{\formatdate{2}{2}{2021}. Redigeret \formatdate{6}{5}{2026}}
\author{Frederik B. Truelsen \& Redigeret af Kristian Anton Hedegaard}

\begin{document}

\maketitle

\section{Lageret}

Lageret skal MINIMUM indeholde følgende antal fustager inden en barvagt.

\begin{itemize}
\item 5 Extra Pilsner (EP)
\item 1 - 2 (billige) Specialøl (???ansvarlige)
\item 1 - 2 (dyre) Specialøl (Specialølsansvarlig)
\item 1 - 2 Cider (Specialølsansvarlig)
\item 2 flasker med CO$_{2}$
\end{itemize}

\subsection{Lagerorden}

Ved gennemgang af lageret, er det en rigtig god ide at flytte eventuelt resterende øl frem tættest på døren.
Dette tjener to formål, vi sikre at bartenderne tager det ældste øl først og giver plads til at Aarhus Bryghus
sætter øllene ind bagerest i lokalet.

\section{Specialtilfælde}

Overstående lagerstatus dækker det almindelig ugentlige behov for at drive Fredagscaféen,
men i tilfælde af andre events med udlejning af vores lille anlæg, specielle events i baren og
i forbindelse med lang barvagt, skal der være ekstra øl.

\subsection{Udlejning}

I tilfælde af en udlejning skal man bestille det ekstra som der er behov for i forbindelse med
udlejningen. I forbindelse store udlejninger som f.eks. Secoya og Institut for Datalogi og andre
festlige lidt større arrangementer  kan det være en god ide at have lidt ekstra fustager i baghånden.

\subsection{Events}

Ved events såsom lange barvagter, sponsorbarer og andre events hvor firmaer udefra kommer forbi skal der forventes
at man skal have et lager på omkring 10 - 20 fustager øl. Dette kommer an på mængden aftalt med eventansvarlige. Man kan følgende uge lave en modsvarende mindre bestilling,
og man behøver ikke bestille øl kan også bestille en tømning af tomme fustager.

\subsubsection{Saint Patrick's Day}

I forbindelse med dette event kan der bestilles øl tilsat grøn farve hvis det bestilles tre uger i forvejen og hvor mange
fustager vi gerne vil have med farve.

\subsection{Ferie}

I sommer- og juleferien er der meget stille og man skal være opmærksom på at en normal bestilling
godt kan holde flere uger, det variere meget gennem ferien og første del af ferien og
midt august ligner mere normalen dog stadig med forbehold for en pæn variation.


\section{Bestilling}

Bestilling foregår primært ved at skrive til vores sælger Christopher på tlf. 27 51 16 55. Alternativt kan man sende en mail til Niels Buchwald på \\
\href{mailto:nfb@aarhusbryghus.dk}{nfb@aarhusbryghus.dk}. Dette må meget gerne gøres senst tirsdag i den uge man har brug for bestillingen.
I beskeden specificerer man hvor mange tomme fustager og CO$_{2}$ flasker vi sender retur, samt hvor meget vi ønsker at bestille,
svarende til at genoprette ovenstående lager status og eventuelt ekstra man har brug i den kommende uge.
Aarhus Bryghus har selv en nøgle og står selv for det praktiske omkring leveringen


\subsection{Holdbarhed}

I forbindelse med store bestillinger og ferier kan det forekomme at nogle af øllene kommer til at løbe over dato.
Aarhus Bryghus kører med måned fra øllen er produceret hvilket også svarer til tæt på den dag de blev leveret.
Vi har ikke haft problemer med at sælge øllen op til en måneds tid efter den dato mærkning, men smag på øllen
om den smager rigtig. Kommer øllen tæt på at gå over den udvidet periode sælges den billigere,
eventuelt omkring købspris afrundet til nærmeste tal deleligt med fem.
Retur foregår ved at man sætter en form for mærkning på dem at de skal retur, og de stilles sammen
med de tomme fustager.



\section{Adgangskort}
Aarhus Bryghus har adgangskort til p-kælderen. Dette udløber pt. hver 3. måned. Derfor er det VIGTIGT at man skriver til Tina (rudolph@cs.au.dk) eller Mette (mmd@cs.au.dk) inden kortet udløber. Alternativt kan cardaccess@cs.au.dk også videresende opgaven. \\
Har man ikke nået at forny kortet inden det udløber, skal man være fysisk til stede til at åbne p-kælderen for Aarhus Bryghus, når de kommer for at hente pant og levere øl. Skriv det til Christopher/Niels ved bestilling, så de ved det, og aftal et tidspunkt hvor I kan mødes. \\
Aarhus Bryghus' kort har følgende oplysninger, som skal sendes med ved fornyelse:
\begin{itemize}
    \item EXTERNAL IDENTITY CARD
    \item Aarhus Bryghus Fredagscafe, øl-depot ADA
    \item Kortnr. 060546
    \item Nummer ved magnetstribe: 106820
\end{itemize}
Der er sat mail-reminder op til best@fredagscafeen.dk cirka en uge før kortet udløber, sådan at man ikke glemmer at forny kortet i tide.


\end{document}
